%% ============================================================================
\section{Proofs}
\label{app:proofs}
%% ============================================================================

Throughout, $\xs$ is a fixed committed world, priors lie in $(0,1)$ and weights in $[0,1)$, so
every world has positive mass under every $\PR$. For a scope $e$ write $V\subseteq\{0,1\}^{|e|}$
for a violated set with $\varnothing\ne V\ne\{0,1\}^{|e|}$; a relation with $V=\varnothing$ or
$V=\{0,1\}^{|e|}$ contributes a constant factor and leaves $\PR$ unchanged.

\subsection{The Score (\S\ref{sec:score})}
\label{app:proof-score}

\paragraph{Lemma~\ref{lem:edit}.} Let $\tilde P(x)=\pi(x)\prod_{r'\in\R}\psi_{r'}(x)$, so
$\PR=\tilde P/Z_\R$. After adding $r$, the unnormalized mass is $\tilde P\psi_r$, with normalizer
$Z_{\R\uplus\{r\}}=\sum_x\tilde P(x)\psi_r(x)=Z_\R\,\E_{\PR}[\psi_r]$. Hence
$P_{\R\uplus\{r\}}(\xs)=\PR(\xs)\psi_r(\xs)/\E_{\PR}[\psi_r]$; taking logarithms and subtracting
the constant $\log\pi(\xs)$ gives the identity. For $\psi_r$ of Eq.~\ref{eq:psi},
$\E_{\PR}[\psi_r]=1-p_rq_r$ with $q_r=\PR(x_{e_r}\in V(\kappa_r))\in(0,1)$, and $\psi_r(\xs)$ is $1$ for a
support and $1-p_r$ for a conflict, which gives the two closed forms and their signs. $\square$

\paragraph{Theorem~\ref{thm:mono}.} By transitivity it suffices to show that each elementary
improvement does not decrease $\Lambda$, and strictly increases it when it changes a weight by a
non-zero amount. Fix the improvement and let $\R_0$ be the graph without the edited relation $r$
(for an addition, the graph before; for a deletion, the graph after). Write
$q=P_{\R_0}(x_e\in V)$; by positivity and the non-triviality of $V$, $0<q<1$. Define
$g(p)=\Lambda(\R_0\uplus\{(e,\kappa,p)\})-\Lambda(\R_0)$. By Lemma~\ref{lem:edit},
$g^+(p)=-\log(1-pq)$ for a support and $g^-(p)=\log(1-p)-\log(1-pq)$ for a conflict, with
$g^\pm(0)=0$, and
\begin{gather*}
\tfrac{d}{dp}g^+=\tfrac{q}{1-pq}>0,\\
\tfrac{d}{dp}g^-=-\tfrac{1-q}{(1-p)(1-pq)}<0\quad\text{on }[0,1).
\end{gather*}
(i) Adding a support of weight $p>0$ changes $\Lambda$ by $g^+(p)>0$. (ii) Raising a support's
weight from $p$ to $p'>p$ changes it by $g^+(p')-g^+(p)>0$. (iii) Deleting a conflict of weight
$p>0$ changes it by $-g^-(p)>0$. (iv) Lowering a conflict's weight from $p$ to $p'<p$ changes it
by $g^-(p')-g^-(p)>0$. A zero-weight addition or deletion leaves $\Lambda$ unchanged. Parallel
copies of the edited relation are part of $\R_0$, so multisets are handled. $\LCS$ is a strictly
increasing function of $\Lambda$ for a fixed claim set and prior, which gives the statement for
$\LCS$. Nothing in the argument uses $|e|=2$, and $\xs$
enters only through which relations are supports, that is, through $\psi_r(\xs)\in\{1,1-p_r\}$.
$\square$

\paragraph{A three-claim example by hand.} Take three claims, all verified true, at $\pi_i=\frac12$ (so $\xs=\one$
under either polarity), with $a_1\Rightarrow a_2$ and $a_2\Rightarrow a_3$ at weight $0.9$ and a
contradiction between $a_1$ and $a_3$ at $0.8$ (Table~\ref{tab:worked}). With only the chain, the
violating worlds are those with $(a_1,a_2)=(1,0)$ or $(a_2,a_3)=(1,0)$; no world violates both,
since that would need $a_2=0$ and $a_2=1$. The unnormalized masses ($\times8$) are
$1,1,0.1,1,0.1,0.1,0.1,1$ for $x=000,\dots,111$, summing to $4.4$, so
$\PR(\one)=1/4.4=\wkChainP$, $\Lambda=\log(8/4.4)=\wkChainLam$, and
$\PR(a_1{=}a_3{=}1)=1.1/4.4=\wkPOneThree$. Adding the contradiction therefore costs exactly
$\log\frac{0.2}{1-0.8\cdot\wkPOneThree}=\wkConflictGain$, as Lemma~\ref{lem:edit} predicts.
Weakening it to $0.44$ recovers most of the loss. The posterior marginals under the chain
are $(0.295,0.500,0.705)$: sign-definite entailments lower the premise's marginal.

\begin{table}[h]
\centering\small
\setlength{\tabcolsep}{5pt}
\begin{tabular}{@{}lccc@{}}
\toprule
Relations & $\PR(\one)$ & $\Lambda$ & $\LCS$ \\
\midrule
none (bag of claims)          & \wkBagP & \wkBagLam & \wkBagLCS \\
contradiction only           & \wkConflictP & \wkConflictLam & \wkConflictLCS \\
chain                        & \wkChainP & \wkChainLam & \wkChainLCS \\
chain + contradiction         & \wkChainConflictP & \wkChainConflictLam & \wkChainConflictLCS \\
chain + weakened contradiction & \wkChainResolvedP & \wkChainResolvedLam & \wkChainResolvedLCS \\
\bottomrule
\end{tabular}
\caption{Three claims, uniform priors, exact values.}
\label{tab:worked}
\end{table}

\subsection{Why This Score (\S\ref{sec:why})}
\label{app:proof-why}

\paragraph{Theorem~\ref{thm:necessity}.} Let $c^\ast=\xs_e$. By Lemma~\ref{lem:edit}, adding
$\psi$ changes $\Lambda$ by $\log\psi(c^\ast)-\log\E_{\PR}[\psi]$, and $\E_{\PR}[\psi]$ is a
convex combination of the values of $\psi$ with weights $\PR(x_e=c)$.
\emph{(If.)} If $\psi(c^\ast)=\max\psi$ the expectation is at most $\psi(c^\ast)$ and the change is
$\ge0$; symmetrically for the minimum.
\emph{(Only if.)} Suppose $\psi(c)>\psi(c^\ast)$ for some cell $c$; we build a relation graph,
with the priors untouched, under which adding $\psi$ lowers $\Lambda$. For every cell $d\ne c$ of
the scope $e=(s,t)$, some coupling has violated set exactly $\{d\}$: $\{(1,0)\}$ is
$a_s\Rightarrow a_t$, $\{(0,1)\}$ is $a_t\Rightarrow a_s$, $\{(1,1)\}$ is \con{} and $\{(0,0)\}$ is
\cnc{}. Let $\R_w$ contain one such relation for each $d\ne c$, all of weight $w$. Every world with
$x_e\ne c$ is multiplied by $1-w$, so $\PR(x_e=c)\to1$ as $w\to1$, whatever the priors. Hence
$\E_{P_{\R_w}}[\psi]\to\psi(c)>\psi(c^\ast)$, and for $w$ close enough to $1$ the change is
negative. The monotone-down case is symmetric. The argument needs a singleton-violation relation for
every cell of the scope, which the pairwise couplings provide; for scopes of higher arity the
theorem holds under the same assumption. \emph{Weight changes.} For a parametric family
$\psi_p$ added to a graph $\R_0$, the change in $\Lambda$ is exactly
$-\log\E_{P_{\R_0}}[\psi_p/\psi_p(c^\ast)]$, so it is monotone in $p$ in every context if each likelihood ratio $\psi_p(c)/\psi_p(c^\ast)$ is
non-increasing in $p$ for supports and non-decreasing for conflicts. This condition is sufficient;
we do not claim it is necessary. The potentials of Eq.~\ref{eq:psi} satisfy it, since the ratios
are $1$, $1-p$ or $1/(1-p)$. $\square$

\paragraph{Proposition~\ref{prop:unique}.} Let $x\ne\xs$, pick $i$ with $x_i\ne\xs_i$ and any
$j\ne i$. As above, every singleton $\{(x_i,x_j)\}$ is the violated set of some coupling on
$(a_i,a_j)$; let $r$ be that relation. Since $(\xs_i,\xs_j)\ne(x_i,x_j)$, $r$ is a support
relative to $\xs$, so adding it with weight $p>0$ is an elementary improvement. It multiplies
$\PR(x)$ by $\psi_r(x)/\E_{\PR}[\psi_r]=(1-p)/(1-pq)$ with $q=\PR(x_e\in V)<1$, which is less
than $1$.
$\square$

\subsection{Further Properties (\S\ref{sec:props})}
\label{app:proof-props}

\paragraph{Proposition~\ref{prop:props}, (i)--(iii).} (i) $P_\varnothing=\pi$. (ii) Build $\R$ by adding its
relations one at a time to $\varnothing$. Each addition of a support raises $\Lambda$ and each
addition of a conflict lowers it (Lemma~\ref{lem:edit}). (iii) With disjoint scopes the
unnormalized measure factorizes as
$\bigl[\pi^{(1)}\prod_{\R_1}\psi\bigr]\bigl[\pi^{(2)}\prod_{\R_2}\psi\bigr]\pi^{(0)}$ over the
claims touched by $\R_1$, those touched by $\R_2$ and the rest. Hence $Z_\R=Z_{\R_1}Z_{\R_2}$, and
the conflict sum in Eq.~\ref{eq:lambda} splits. $\square$

\paragraph{Proposition~\ref{prop:props}(iv).} The map $x\mapsto x'$ that flips coordinate $i$ is a
bijection of $\{0,1\}^n$. It sends $\xs$ to the new committed world, $\pi(x)$ to the new prior of
$x'$, since $\pi_i^{x_i}(1-\pi_i)^{1-x_i}=(1-\pi_i)^{x'_i}\pi_i^{1-x'_i}$, and each $\psi_r(x)$ to
the rewritten relation's $\psi_{r'}(x')$, since negating an endpoint maps the violated set
cell-wise. Hence $Z_\R$, $\PR(\xs)$ and $\pi(\xs)$ are preserved, and so are $\Lambda$ and $\LCS$;
$\xs_{e_r}\in V(\kappa_r)$ iff ${\xs}'_{e_{r'}}\in V(\kappa_{r'})$, so signs are preserved. The rewritten
violated set is again a coupling of Def.~\ref{def:coupling} (App.~\ref{app:framework}). $\square$

\paragraph{Proposition~\ref{prop:props}, (v)--(vii).} (v) The upper bound is $\PR(\xs)\le1$. The lower bound
follows from $Z_\R=\E_\pi[\prod\psi]\le1$, since every potential is at most $1$. (vi) For a
support, $V\not\ni\xs_e$, so $q_r=\PR(x_e\in V)\le1-\PR(\xs)$, and $-\log(1-p_rq_r)$ is
increasing in $q_r$. (vii) After adding a support $r$ with violation probability $q$, the new
violation probability is $q'=q(1-p)/(1-pq)<q$ for $p>0$ and $q<1$, so a second copy gains
$-\log(1-pq')<-\log(1-pq)$. $\square$

\paragraph{Saturation in numbers.} Successive copies of one equivalence ($p=0.9$) on two fair
claims gain $\satGainOne$, $\satGainTwo$, $\satGainThree$, $\satGainFour$, \dots{} nats; after five
copies $\Lambda=\satLamFive$, and it converges to $\log2$ rather than to the bound $\satBound$: an
equivalence cannot distinguish ``both hold'' from ``both fail''. Under the prior-expectation score
of App.~\ref{app:earlier}, one to six copies drive the score from $\padBalOne$ to $\padBalSix$,
while $\LCS$ goes from $\padLCSOne$ to $\padLCSSix$ and does not move, to three decimals, after the
third copy.

